\chapter{Conclusion}
\label{chap:conclusion}

\section{Summary}

This report has proposed BiasAperture, a diagnostic and evaluative software platform for auditing demographic accuracy disparities in third-party facial analysis models, and specified it in enough detail to be built against by a two-person team over an eight-week schedule. \Cref{chap:litreview} showed that a mature fairness-metrics and mitigation-techniques literature exists without a corresponding reusable auditing artefact, at precisely the moment the \gls{eu} \gls{ai} Act's Article~10 obligations make that gap a compliance question rather than only an academic one. \Cref{chap:requirements,chap:methodology} then specified how the platform closes it, against each of \cref{chap:intro}'s five specific objectives in turn: a modular architecture built around FR-001 and FR-002 satisfies the first; AIF360 and Fairlearn as independent, cross-validating backends for the four Core~Four metrics, with chi-squared and bootstrap significance testing, satisfy the second; a report structured against the Model Cards and Datasheets conventions surveyed in \cref{chap:litreview} satisfies the third; validation against the current case study, which uses FairFace since UTKFace was profiled and cut from the implementation scope, with a FairFace-trained \gls{cnn} baseline as the minimum case study, satisfies the fourth; and the Article~10 and \gls{nist} \gls{ai} \gls{rmf} mapping detailed in \cref{sec:regmapping} satisfies the fifth.

Consistent with the scope stated in \cref{chap:intro}, BiasAperture remains strictly diagnostic throughout this design: it reports where a disparity exists and how statistically strong it is, and it does not mitigate bias, retrain a model, or generate synthetic demographic data. \Cref{sec:cutlist}'s descoping order exists precisely so that this diagnostic core, ingestion, one model interface, the fairness engine, one report format, and the scope-boundary statement itself, is the one part of the proposal that is never at risk under timeline pressure, whatever else is cut first.

\section{Future Work}

The descoping order in \cref{tab:cutlist} doubles as a map of the platform's nearest extension points once the eight-week schedule in \cref{sec:schedule} is complete: a web \gls{ui} over the existing \gls{cli}, \gls{pdf} export alongside the \gls{html} report, direct in-process inference restored alongside predictions-file ingestion, and UTKFace (currently cut per Cut-List \#2) brought back in as a fully supported secondary benchmark once its label-noise limitation, discussed in \cref{chap:litreview}, is more thoroughly characterised. Beyond the current cut-list, the clearest longer-term direction follows from the diagnostic and mitigation split Dehdashtian et al.'s survey draws \cite{dehdashtian2024survey}: BiasAperture is deliberately positioned on the diagnostic side of that split, and a natural next stage of the same research programme is a mitigation-facing companion that consumes BiasAperture's findings as input rather than folding mitigation into the diagnostic platform itself, keeping the separation of concerns this proposal establishes intact rather than eroding it.
